\documentclass[11pt,a4paper]{article}

% Paquetes
\usepackage[utf8]{inputenc}
\usepackage[spanish]{babel}
\usepackage{geometry}
\usepackage{listings}
\usepackage{xcolor}
\usepackage{hyperref}
\usepackage{graphicx}
\usepackage{fancyhdr}
\usepackage{tcolorbox}
\usepackage{enumitem}
\usepackage{titlesec}

% Configuración de página
\geometry{
    a4paper,
    left=2.5cm,
    right=2.5cm,
    top=3cm,
    bottom=3cm
}

% Configuración de encabezados y pies de página
\pagestyle{fancy}
\fancyhf{}
\fancyhead[L]{Guía Completa de Nginx}
\fancyhead[R]{\thepage}
\renewcommand{\headrulewidth}{0.4pt}

% Configuración de colores
\definecolor{codebackground}{RGB}{248,248,248}
\definecolor{codekeyword}{RGB}{0,0,255}
\definecolor{codecomment}{RGB}{0,128,0}
\definecolor{codestring}{RGB}{163,21,21}

% Configuración de listings para código
\lstdefinestyle{nginx}{
    language=bash,
    backgroundcolor=\color{codebackground},
    commentstyle=\color{codecomment},
    keywordstyle=\color{codekeyword},
    stringstyle=\color{codestring},
    basicstyle=\ttfamily\footnotesize,
    breakatwhitespace=false,
    breaklines=true,
    captionpos=b,
    keepspaces=true,
    numbers=left,
    numbersep=5pt,
    showspaces=false,
    showstringspaces=false,
    showtabs=false,
    tabsize=2,
    frame=single,
    rulecolor=\color{black!30}
}

\lstset{style=nginx}

% Configuración de hyperref
\hypersetup{
    colorlinks=true,
    linkcolor=blue,
    filecolor=magenta,
    urlcolor=cyan,
    pdftitle={Guía Completa de Nginx},
    pdfauthor={AlexSilvaa9},
}

% Título del documento
\title{
    \Huge\textbf{Guía Completa de Nginx} \\
    \Large Todo lo que necesitas saber sobre Nginx y Virtual Hosts
}
\author{AlexSilvaa9}
\date{\today}

\begin{document}

\maketitle
\newpage

\tableofcontents
\newpage

% ===================================
% SECCIÓN 1: INTRODUCCIÓN
% ===================================
\section{Introducción a Nginx}

Nginx (pronunciado "engine-x") es un servidor web de código abierto, conocido por su alto rendimiento, estabilidad, configuración simple y bajo consumo de recursos.

\subsection{Características principales}

\begin{itemize}[leftmargin=*]
    \item Servidor web HTTP de alto rendimiento
    \item Proxy reverso para HTTP, HTTPS, SMTP, POP3 e IMAP
    \item Balanceador de carga
    \item Caché HTTP
    \item Manejo eficiente de archivos estáticos
    \item Soporte para FastCGI, uWSGI, SCGI
\end{itemize}

% ===================================
% SECCIÓN 2: VIRTUAL HOSTS
% ===================================
\section{Virtual Hosts (vhosts)}

\subsection{¿Qué es un Virtual Host?}

Un \textbf{virtual host} es una configuración que permite que un \textbf{mismo servidor web} (una sola máquina con una IP) sirva \textbf{múltiples sitios web o aplicaciones diferentes}.

\subsubsection{Analogía}
Imagina un edificio de apartamentos (el servidor):
\begin{itemize}
    \item \textbf{Una dirección} (una IP)
    \item \textbf{Múltiples apartamentos} (diferentes sitios web)
    \item \textbf{El portero} (Nginx) dirige a cada visitante al apartamento correcto según el nombre que pida
\end{itemize}

\subsection{Tipos de Virtual Hosts}

\subsubsection{Basados en nombre (Name-based)}
El más común. Usa el \textbf{nombre de dominio} para distinguir sitios:
\begin{itemize}
    \item \texttt{ejemplo1.com} → Sitio A
    \item \texttt{ejemplo2.com} → Sitio B
    \item \textbf{Misma IP}, diferentes dominios
\end{itemize}

\subsubsection{Basados en IP (IP-based)}
Cada sitio tiene su propia IP:
\begin{itemize}
    \item \texttt{192.168.1.10} → Sitio A
    \item \texttt{192.168.1.11} → Sitio B
\end{itemize}

\subsubsection{Basados en puerto (Port-based)}
Mismo dominio/IP, diferentes puertos:
\begin{itemize}
    \item \texttt{ejemplo.com:80} → Sitio A
    \item \texttt{ejemplo.com:8080} → Sitio B
\end{itemize}

% ===================================
% SECCIÓN 3: ARQUITECTURA
% ===================================
\section{Arquitectura de Nginx}

\subsection{Modelo de procesos}

Nginx usa un modelo \textbf{asíncrono y basado en eventos} (event-driven):

\begin{tcolorbox}[colback=blue!5!white,colframe=blue!75!black,title=Estructura de procesos]
\begin{verbatim}
            ┌─────────────────┐
            │  Master Process │  (root, gestiona workers)
            └────────┬────────┘
                     │
    ┌────────────────┼────────────────┐
    │                │                │
┌───▼────┐      ┌────▼────┐      ┌───▼────┐
│ Worker │      │ Worker  │      │ Worker │
│Process │      │ Process │      │Process │
└────────┘      └─────────┘      └────────┘
(atiende miles de conexiones cada uno)
\end{verbatim}
\end{tcolorbox}

\subsection{Ventajas del modelo de Nginx}
\begin{itemize}
    \item Un worker puede manejar miles de conexiones simultáneas
    \item Muy eficiente en memoria
    \item No bloquea operaciones (non-blocking I/O)
\end{itemize}

\subsection{Comparación con Apache}
\begin{itemize}
    \item \textbf{Apache}: Crea un proceso/thread por conexión (más pesado)
    \item \textbf{Nginx}: Event-driven, maneja múltiples conexiones por worker
\end{itemize}

% ===================================
% SECCIÓN 4: ESTRUCTURA DE CONFIGURACIÓN
% ===================================
\section{Estructura de Configuración}

\subsection{Directorios principales}

\begin{lstlisting}[language=bash, caption=Estructura de directorios de Nginx]
/etc/nginx/
├── nginx.conf              # Configuración principal
├── sites-available/        # Configuraciones disponibles
│   ├── sitio1.conf
│   ├── sitio2.conf
│   └── default
├── sites-enabled/          # Symlinks a sitios activos
│   ├── sitio1.conf -> ../sites-available/sitio1.conf
│   └── default -> ../sites-available/default
├── conf.d/                 # Configuraciones adicionales
├── snippets/               # Fragmentos reutilizables
└── mime.types              # Tipos MIME
\end{lstlisting}

\subsection{Niveles de contexto}

\begin{lstlisting}[caption=Estructura completa de nginx.conf]
# CONTEXTO MAIN (global)
user www-data;
worker_processes auto;
pid /run/nginx.pid;

events {
    # CONTEXTO EVENTS
    worker_connections 1024;
    use epoll;
}

http {
    # CONTEXTO HTTP
    include /etc/nginx/mime.types;
    default_type application/octet-stream;
    
    sendfile on;
    keepalive_timeout 65;
    
    access_log /var/log/nginx/access.log;
    error_log /var/log/nginx/error.log;
    
    upstream backend {
        server localhost:3000;
        server localhost:3001;
    }
    
    server {
        # CONTEXTO SERVER (virtual host)
        listen 80;
        server_name ejemplo.com;
        
        location / {
            # CONTEXTO LOCATION
            root /var/www/html;
            index index.html;
        }
    }
    
    include /etc/nginx/conf.d/*.conf;
    include /etc/nginx/sites-enabled/*;
}
\end{lstlisting}

% ===================================
% SECCIÓN 5: EJEMPLOS DE VIRTUAL HOSTS
% ===================================
\section{Ejemplos de Virtual Hosts}

\subsection{Sitio web estático simple}

\begin{lstlisting}[caption=sites-available/miapp.com.conf]
server {
    listen 80;
    listen [::]:80;
    
    server_name miapp.com www.miapp.com;
    
    root /var/www/miapp;
    index index.html index.htm;
    
    access_log /var/log/nginx/miapp-access.log;
    error_log /var/log/nginx/miapp-error.log;
    
    location / {
        try_files $uri $uri/ =404;
    }
}
\end{lstlisting}

\subsection{Proxy reverso a aplicación}

\begin{lstlisting}[caption=Proxy a Node.js/Python/etc]
server {
    listen 80;
    server_name api.miapp.com;
    
    location / {
        proxy_pass http://localhost:3000;
        proxy_http_version 1.1;
        
        proxy_set_header Upgrade $http_upgrade;
        proxy_set_header Connection 'upgrade';
        proxy_set_header Host $host;
        proxy_set_header X-Real-IP $remote_addr;
        proxy_set_header X-Forwarded-For $proxy_add_x_forwarded_for;
        proxy_set_header X-Forwarded-Proto $scheme;
        
        proxy_cache_bypass $http_upgrade;
    }
}
\end{lstlisting}

\subsection{HTTPS con SSL/TLS}

\begin{lstlisting}[caption=Configuración SSL completa]
server {
    listen 80;
    server_name miapp.com www.miapp.com;
    return 301 https://$server_name$request_uri;
}

server {
    listen 443 ssl http2;
    listen [::]:443 ssl http2;
    
    server_name miapp.com www.miapp.com;
    
    ssl_certificate /etc/ssl/certs/miapp.crt;
    ssl_certificate_key /etc/ssl/private/miapp.key;
    
    ssl_protocols TLSv1.2 TLSv1.3;
    ssl_ciphers HIGH:!aNULL:!MD5;
    ssl_prefer_server_ciphers on;
    
    root /var/www/miapp;
    index index.html;
    
    location / {
        try_files $uri $uri/ =404;
    }
}
\end{lstlisting}

\subsection{Múltiples aplicaciones en subdominios}

\begin{lstlisting}[caption=Multi-sitio con subdominios]
# Aplicación principal
server {
    listen 80;
    server_name miempresa.com www.miempresa.com;
    root /var/www/principal;
    index index.html;
}

# Blog
server {
    listen 80;
    server_name blog.miempresa.com;
    
    location / {
        proxy_pass http://localhost:2368;
        proxy_set_header Host $host;
    }
}

# API
server {
    listen 80;
    server_name api.miempresa.com;
    
    location / {
        proxy_pass http://localhost:3000;
        proxy_set_header Host $host;
    }
}
\end{lstlisting}

% ===================================
% SECCIÓN 6: TRY_FILES
% ===================================
\section{La directiva try\_files}

\subsection{¿Qué es try\_files?}

\texttt{try\_files} es una directiva que \textbf{intenta servir archivos en un orden específico} hasta encontrar uno que exista.

\subsection{Sintaxis}
\begin{lstlisting}
try_files archivo1 archivo2 ... fallback;
\end{lstlisting}

\subsection{Funcionamiento paso a paso}

Cuando llega una petición \texttt{http://miapp.com/about}:

\begin{enumerate}
    \item \textbf{Paso 1}: Nginx busca \texttt{\$uri} → \texttt{/var/www/miapp/about} (archivo)
    \item \textbf{Paso 2}: Nginx busca \texttt{\$uri/} → \texttt{/var/www/miapp/about/} (directorio)
    \item \textbf{Paso 3}: Nginx sirve el fallback → \texttt{/index.html}
\end{enumerate}

\subsection{Casos de uso}

\subsubsection{SPA (React, Vue, Angular)}
\begin{lstlisting}
location / {
    try_files $uri $uri/ /index.html;
}
\end{lstlisting}

\subsubsection{Sitio estático}
\begin{lstlisting}
location / {
    try_files $uri $uri/ =404;
}
\end{lstlisting}

\subsubsection{WordPress/PHP}
\begin{lstlisting}
location / {
    try_files $uri $uri/ /index.php?$args;
}
\end{lstlisting}

% ===================================
% SECCIÓN 7: DIRECTIVAS IMPORTANTES
% ===================================
\section{Directivas Importantes}

\subsection{Listen}
\begin{lstlisting}
listen 80;                    # IPv4, puerto 80
listen [::]:80;               # IPv6, puerto 80
listen 443 ssl;               # HTTPS
listen 8080;                  # Puerto personalizado
\end{lstlisting}

\subsection{Server Name}
\begin{lstlisting}
server_name ejemplo.com;
server_name ejemplo.com www.ejemplo.com;
server_name *.ejemplo.com;
server_name ~^(?<subdomain>.+)\.ejemplo\.com$;
\end{lstlisting}

\subsection{Root e Index}
\begin{lstlisting}
root /var/www/html;
index index.html index.php;
\end{lstlisting}

\subsection{Location - Prioridades}

Las directivas \texttt{location} se evalúan en este orden:

\begin{enumerate}
    \item \texttt{location = /exacto} - Match exacto (PRIORIDAD MÁXIMA)
    \item \texttt{location \^{}\~{} /img/} - Preferential prefix
    \item \texttt{location \~{} \.php\$} - Regex case-sensitive
    \item \texttt{location \~{}* \.(jpg|png)\$} - Regex case-insensitive
    \item \texttt{location /} - Prefix match (PRIORIDAD MÍNIMA)
\end{enumerate}

\begin{lstlisting}[caption=Ejemplo de prioridades]
server {
    listen 80;
    server_name test.local;
    
    location = / {
        return 200 "Exacto: raiz\n";
    }
    
    location / {
        return 200 "Prefix: cualquier cosa\n";
    }
    
    location /imagenes/ {
        return 200 "Prefix: imagenes\n";
    }
    
    location ~ \.png$ {
        return 200 "Regex: archivos PNG\n";
    }
}
\end{lstlisting}

% ===================================
% SECCIÓN 8: DNS Y VIRTUAL HOSTS
% ===================================
\section{DNS y Virtual Hosts}

\subsection{Flujo completo}

\subsubsection{Paso 1: Configuración DNS}

En tu proveedor DNS (Cloudflare, GoDaddy, Route53):

\begin{tcolorbox}[colback=green!5!white,colframe=green!75!black,title=Registros DNS]
\begin{verbatim}
Type    Name              Value           TTL
A       miapp.com         203.0.113.10    3600
A       blog.com          203.0.113.10    3600
A       tienda.com        203.0.113.10    3600
\end{verbatim}
\end{tcolorbox}

✅ \textbf{Múltiples dominios apuntan a la MISMA IP}

\subsubsection{Paso 2: El navegador hace la petición}

\begin{lstlisting}[language=html]
GET / HTTP/1.1
Host: miapp.com
\end{lstlisting}

\textbf{Clave:} El navegador envía el nombre del dominio en la cabecera \texttt{Host:}

\subsubsection{Paso 3: Nginx elige el virtual host}

Nginx compara la cabecera \texttt{Host:} con \texttt{server\_name}:

\begin{itemize}
    \item Si \texttt{Host: miapp.com} → sirve \texttt{/var/www/miapp}
    \item Si \texttt{Host: blog.com} → sirve \texttt{/var/www/blog}
    \item Si \texttt{Host: tienda.com} → sirve \texttt{/var/www/tienda}
\end{itemize}

% ===================================
% SECCIÓN 9: HABILITAR SITIOS
% ===================================
\section{Habilitar y Deshabilitar Sitios}

\subsection{Concepto}

Nginx usa dos directorios:
\begin{itemize}
    \item \texttt{sites-available/} - Todas las configuraciones (activas o no)
    \item \texttt{sites-enabled/} - Solo las configuraciones ACTIVAS
\end{itemize}

\subsection{Proceso completo}

\subsubsection{1. Crear configuración}
\begin{lstlisting}[language=bash]
sudo nano /etc/nginx/sites-available/miapp.conf
\end{lstlisting}

\subsubsection{2. Habilitar el sitio (crear symlink)}
\begin{lstlisting}[language=bash]
sudo ln -s /etc/nginx/sites-available/miapp.conf \
           /etc/nginx/sites-enabled/
\end{lstlisting}

\subsubsection{3. Verificar y recargar}
\begin{lstlisting}[language=bash]
sudo nginx -t
sudo nginx -s reload
\end{lstlisting}

\subsubsection{4. Deshabilitar sitio}
\begin{lstlisting}[language=bash]
sudo rm /etc/nginx/sites-enabled/miapp.conf
sudo nginx -s reload
\end{lstlisting}

\textbf{Importante:} Solo se borra el symlink, la configuración permanece en \texttt{sites-available/}

% ===================================
% SECCIÓN 10: PROXY REVERSO
% ===================================
\section{Proxy Reverso}

\subsection{Headers esenciales}

\begin{lstlisting}[caption=Configuración completa de proxy]
location / {
    proxy_pass http://backend;
    
    # Headers esenciales
    proxy_set_header Host $host;
    proxy_set_header X-Real-IP $remote_addr;
    proxy_set_header X-Forwarded-For $proxy_add_x_forwarded_for;
    proxy_set_header X-Forwarded-Proto $scheme;
    proxy_set_header X-Forwarded-Host $host;
    proxy_set_header X-Forwarded-Port $server_port;
    
    # WebSockets
    proxy_http_version 1.1;
    proxy_set_header Upgrade $http_upgrade;
    proxy_set_header Connection "upgrade";
    
    # Timeouts
    proxy_connect_timeout 60s;
    proxy_send_timeout 60s;
    proxy_read_timeout 60s;
    
    # Buffers
    proxy_buffering on;
    proxy_buffer_size 4k;
    proxy_buffers 8 4k;
    
    proxy_cache_bypass $http_upgrade;
}
\end{lstlisting}

\subsection{Upstream y Load Balancing}

\subsubsection{Round-robin (por defecto)}
\begin{lstlisting}
upstream backend {
    server backend1.com;
    server backend2.com;
    server backend3.com;
}
\end{lstlisting}

\subsubsection{Least connections}
\begin{lstlisting}
upstream backend {
    least_conn;
    server backend1.com;
    server backend2.com;
}
\end{lstlisting}

\subsubsection{IP Hash}
\begin{lstlisting}
upstream backend {
    ip_hash;
    server backend1.com;
    server backend2.com;
}
\end{lstlisting}

\subsubsection{Con pesos y parámetros}
\begin{lstlisting}
upstream backend {
    server backend1.com weight=3;
    server backend2.com max_fails=3 fail_timeout=30s;
    server backend3.com backup;
    server backend4.com down;
}
\end{lstlisting}

% ===================================
% SECCIÓN 11: CACHÉ
% ===================================
\section{Caché en Nginx}

\subsection{Caché de Proxy}

\begin{lstlisting}[caption=Configuración de caché]
http {
    proxy_cache_path /var/cache/nginx 
                     levels=1:2 
                     keys_zone=my_cache:10m 
                     max_size=1g 
                     inactive=60m 
                     use_temp_path=off;
    
    server {
        location / {
            proxy_pass http://backend;
            proxy_cache my_cache;
            proxy_cache_valid 200 10m;
            proxy_cache_valid 404 1m;
            
            add_header X-Cache-Status $upstream_cache_status;
            
            proxy_cache_key "$scheme$request_method$host$request_uri";
            proxy_cache_bypass $cookie_nocache $arg_nocache;
        }
    }
}
\end{lstlisting}

\subsection{Caché de archivos estáticos}

\begin{lstlisting}
location ~* \.(jpg|jpeg|png|gif|ico|css|js)$ {
    expires 1y;
    add_header Cache-Control "public, immutable";
    access_log off;
}
\end{lstlisting}

% ===================================
% SECCIÓN 12: SEGURIDAD
% ===================================
\section{Seguridad}

\subsection{Rate Limiting}

\begin{lstlisting}[caption=Protección contra DDoS]
http {
    limit_req_zone $binary_remote_addr zone=mylimit:10m rate=10r/s;
    
    server {
        location /api/ {
            limit_req zone=mylimit burst=20 nodelay;
        }
    }
}
\end{lstlisting}

\subsection{Limit de conexiones}

\begin{lstlisting}
http {
    limit_conn_zone $binary_remote_addr zone=addr:10m;
    
    server {
        location /download/ {
            limit_conn addr 2;
        }
    }
}
\end{lstlisting}

\subsection{Control de acceso}

\begin{lstlisting}
location /admin {
    allow 192.168.1.0/24;
    allow 10.0.0.1;
    deny all;
}
\end{lstlisting}

\subsection{Headers de seguridad}

\begin{lstlisting}
add_header X-Frame-Options "SAMEORIGIN" always;
add_header X-Content-Type-Options "nosniff" always;
add_header X-XSS-Protection "1; mode=block" always;
add_header Referrer-Policy "no-referrer-when-downgrade" always;
add_header Content-Security-Policy 
    "default-src 'self' http: https: data: blob: 'unsafe-inline'" always;

server_tokens off;
\end{lstlisting}

\subsection{Autenticación básica}

\begin{lstlisting}[language=bash]
# Crear archivo de contraseñas
sudo htpasswd -c /etc/nginx/.htpasswd usuario1
\end{lstlisting}

\begin{lstlisting}
location /admin {
    auth_basic "Area restringida";
    auth_basic_user_file /etc/nginx/.htpasswd;
}
\end{lstlisting}

% ===================================
% SECCIÓN 13: REDIRECTS Y REWRITES
% ===================================
\section{Redirects y Rewrites}

\subsection{Return (más eficiente)}

\begin{lstlisting}
# Redirect permanente
server {
    listen 80;
    server_name old-domain.com;
    return 301 https://new-domain.com$request_uri;
}

# HTTP a HTTPS
server {
    listen 80;
    return 301 https://$server_name$request_uri;
}

# Redirect de página
location /old-page {
    return 301 /new-page;
}
\end{lstlisting}

\subsection{Rewrite}

\begin{lstlisting}
# Rewrite simple
rewrite ^/old/(.*)$ /new/$1 permanent;

# Con condiciones
if ($host != 'www.ejemplo.com') {
    rewrite ^/(.*)$ https://www.ejemplo.com/$1 permanent;
}
\end{lstlisting}

\textbf{Flags de rewrite:}
\begin{itemize}
    \item \texttt{last} - Detiene y busca nueva location
    \item \texttt{break} - Detiene en esta location
    \item \texttt{redirect} - Redirect temporal (302)
    \item \texttt{permanent} - Redirect permanente (301)
\end{itemize}

% ===================================
% SECCIÓN 14: SSL/TLS
% ===================================
\section{SSL/TLS}

\subsection{Configuración SSL completa}

\begin{lstlisting}[caption=SSL/TLS moderno y seguro]
server {
    listen 443 ssl http2;
    listen [::]:443 ssl http2;
    server_name ejemplo.com;
    
    # Certificados
    ssl_certificate /etc/letsencrypt/live/ejemplo.com/fullchain.pem;
    ssl_certificate_key /etc/letsencrypt/live/ejemplo.com/privkey.pem;
    
    # Protocolos modernos
    ssl_protocols TLSv1.2 TLSv1.3;
    
    # Cifrados seguros
    ssl_ciphers 'ECDHE-ECDSA-AES128-GCM-SHA256:ECDHE-RSA-AES128-GCM-SHA256';
    ssl_prefer_server_ciphers off;
    
    # OCSP Stapling
    ssl_stapling on;
    ssl_stapling_verify on;
    resolver 8.8.8.8 8.8.4.4 valid=300s;
    
    # Sesiones SSL
    ssl_session_cache shared:SSL:10m;
    ssl_session_timeout 10m;
    ssl_session_tickets off;
    
    # HSTS
    add_header Strict-Transport-Security 
        "max-age=31536000; includeSubDomains; preload" always;
    
    # Diffie-Hellman
    ssl_dhparam /etc/nginx/dhparam.pem;
}
\end{lstlisting}

\subsection{Let's Encrypt}

\begin{lstlisting}[language=bash]
# Instalar certbot
sudo apt install certbot python3-certbot-nginx

# Obtener certificado
sudo certbot --nginx -d ejemplo.com -d www.ejemplo.com

# Auto-renovación
sudo certbot renew --dry-run
\end{lstlisting}

% ===================================
% SECCIÓN 15: LOGS
% ===================================
\section{Logs y Monitoring}

\subsection{Formatos de log personalizados}

\begin{lstlisting}
http {
    log_format main '$remote_addr - $remote_user [$time_local] '
                    '"$request" $status $body_bytes_sent '
                    '"$http_referer" "$http_user_agent" '
                    '$request_time';
    
    log_format json escape=json '{'
        '"time":"$time_local",'
        '"remote_addr":"$remote_addr",'
        '"request":"$request",'
        '"status":$status,'
        '"body_bytes_sent":$body_bytes_sent,'
        '"request_time":$request_time'
    '}';
    
    access_log /var/log/nginx/access.log main;
    
    server {
        access_log /var/log/nginx/miapp-access.log json;
        error_log /var/log/nginx/miapp-error.log warn;
    }
}
\end{lstlisting}

\subsection{Analizar logs}

\begin{lstlisting}[language=bash]
# Ver últimas peticiones
tail -f /var/log/nginx/access.log

# IPs con más peticiones
awk '{print $1}' /var/log/nginx/access.log | \
    sort | uniq -c | sort -rn | head -10

# Códigos de estado
awk '{print $9}' /var/log/nginx/access.log | \
    sort | uniq -c | sort -rn

# URLs más visitadas
awk '{print $7}' /var/log/nginx/access.log | \
    sort | uniq -c | sort -rn | head -10
\end{lstlisting}

% ===================================
% SECCIÓN 16: VARIABLES
% ===================================
\section{Variables de Nginx}

\subsection{Variables principales}

\begin{table}[h]
\centering
\begin{tabular}{|l|p{8cm}|}
\hline
\textbf{Variable} & \textbf{Descripción} \\
\hline
\texttt{\$host} & Nombre del host (ejemplo.com) \\
\texttt{\$uri} & URI de la petición (/path/to/file) \\
\texttt{\$request\_uri} & URI completa con query string \\
\texttt{\$args} & Query string (?id=123\&name=test) \\
\texttt{\$remote\_addr} & IP del cliente \\
\texttt{\$server\_addr} & IP del servidor \\
\texttt{\$server\_port} & Puerto del servidor \\
\texttt{\$scheme} & http o https \\
\texttt{\$request\_method} & GET, POST, PUT, etc. \\
\texttt{\$http\_user\_agent} & User agent del navegador \\
\texttt{\$http\_referer} & Página de referencia \\
\texttt{\$status} & Código de estado HTTP \\
\texttt{\$body\_bytes\_sent} & Bytes enviados (sin headers) \\
\texttt{\$request\_time} & Tiempo total de procesamiento \\
\hline
\end{tabular}
\caption{Variables principales de Nginx}
\end{table}

% ===================================
% SECCIÓN 17: PERFORMANCE
% ===================================
\section{Performance Tuning}

\begin{lstlisting}[caption=Optimización de rendimiento]
user www-data;
worker_processes auto;
worker_rlimit_nofile 65535;

events {
    worker_connections 4096;
    use epoll;
    multi_accept on;
}

http {
    # Sendfile y TCP
    sendfile on;
    tcp_nopush on;
    tcp_nodelay on;
    
    # Timeouts
    keepalive_timeout 30;
    keepalive_requests 100;
    reset_timedout_connection on;
    client_body_timeout 10;
    send_timeout 10;
    
    # Buffers
    client_body_buffer_size 128k;
    client_max_body_size 10m;
    client_header_buffer_size 1k;
    large_client_header_buffers 4 4k;
    
    # Gzip
    gzip on;
    gzip_vary on;
    gzip_comp_level 5;
    gzip_min_length 256;
    gzip_types
        text/plain
        text/css
        text/xml
        text/javascript
        application/json
        application/javascript
        application/xml+rss;
    
    # File cache
    open_file_cache max=200000 inactive=20s;
    open_file_cache_valid 30s;
    open_file_cache_min_uses 2;
    open_file_cache_errors on;
}
\end{lstlisting}

% ===================================
% SECCIÓN 18: TROUBLESHOOTING
% ===================================
\section{Troubleshooting}

\subsection{Comandos de diagnóstico}

\begin{lstlisting}[language=bash]
# Verificar sintaxis
sudo nginx -t

# Ver configuración completa
sudo nginx -T

# Ver versión y módulos
nginx -V

# Ver procesos
ps aux | grep nginx

# Ver puertos
netstat -tlnp | grep nginx

# Ver logs en tiempo real
tail -f /var/log/nginx/error.log
\end{lstlisting}

\subsection{Errores comunes y soluciones}

\subsubsection{Could not build server\_names\_hash}
\begin{lstlisting}
http {
    server_names_hash_bucket_size 64;
}
\end{lstlisting}

\subsubsection{413 Request Entity Too Large}
\begin{lstlisting}
http {
    client_max_body_size 100M;
}
\end{lstlisting}

\subsubsection{502 Bad Gateway}
\begin{itemize}
    \item Backend no está corriendo
    \item Firewall bloqueando conexión
    \item Socket PHP-FPM incorrecto
\end{itemize}

\begin{lstlisting}[language=bash]
# Verificar backend
curl http://localhost:3000

# Ver logs
tail -f /var/log/nginx/error.log
\end{lstlisting}

\subsubsection{Permission denied}
\begin{lstlisting}[language=bash]
# Verificar permisos
ls -la /var/www/miapp

# Cambiar dueño
sudo chown -R www-data:www-data /var/www/miapp
\end{lstlisting}

% ===================================
% SECCIÓN 19: COMANDOS ESENCIALES
% ===================================
\section{Comandos Esenciales}

\subsection{Gestión del servicio}

\begin{lstlisting}[language=bash]
# Iniciar
sudo systemctl start nginx

# Detener
sudo systemctl stop nginx

# Reiniciar
sudo systemctl restart nginx

# Recargar (sin downtime)
sudo systemctl reload nginx

# Ver estado
sudo systemctl status nginx

# Auto-start en boot
sudo systemctl enable nginx
\end{lstlisting}

\subsection{Comandos de Nginx}

\begin{lstlisting}[language=bash]
# Test de configuración
sudo nginx -t

# Dump de configuración
sudo nginx -T

# Reload
sudo nginx -s reload

# Stop
sudo nginx -s stop

# Graceful shutdown
sudo nginx -s quit

# Reabrir logs
sudo nginx -s reopen
\end{lstlisting}

% ===================================
% SECCIÓN 20: CASOS DE USO
% ===================================
\section{Casos de Uso Avanzados}

\subsection{Microservicios}

\begin{lstlisting}
server {
    listen 80;
    server_name api.ejemplo.com;
    
    location /users {
        proxy_pass http://localhost:3001;
    }
    
    location /products {
        proxy_pass http://localhost:3002;
    }
    
    location /payments {
        proxy_pass http://localhost:3003;
    }
}
\end{lstlisting}

\subsection{SPA con API backend}

\begin{lstlisting}
server {
    listen 80;
    server_name app.ejemplo.com;
    root /var/www/frontend/dist;
    
    # Frontend
    location / {
        try_files $uri $uri/ /index.html;
    }
    
    # API
    location /api {
        proxy_pass http://localhost:3000;
        proxy_set_header Host $host;
    }
    
    # WebSocket
    location /ws {
        proxy_pass http://localhost:3000;
        proxy_http_version 1.1;
        proxy_set_header Upgrade $http_upgrade;
        proxy_set_header Connection "Upgrade";
    }
    
    # Archivos estáticos
    location /static {
        expires 1y;
        add_header Cache-Control "public, immutable";
    }
}
\end{lstlisting}

\subsection{Multi-tenant (SaaS)}

\begin{lstlisting}
map $http_host $tenant {
    cliente1.miapp.com "cliente1";
    cliente2.miapp.com "cliente2";
    default "demo";
}

server {
    listen 80;
    server_name *.miapp.com;
    
    location / {
        proxy_pass http://localhost:3000;
        proxy_set_header X-Tenant-ID $tenant;
        proxy_set_header Host $host;
    }
}
\end{lstlisting}

% ===================================
% SECCIÓN 21: SNIPPETS REUTILIZABLES
% ===================================
\section{Snippets Reutilizables}

\subsection{SSL Parameters}

\begin{lstlisting}[caption=/etc/nginx/snippets/ssl-params.conf]
ssl_protocols TLSv1.2 TLSv1.3;
ssl_ciphers 'ECDHE-ECDSA-AES128-GCM-SHA256:ECDHE-RSA-AES128-GCM-SHA256';
ssl_prefer_server_ciphers off;
ssl_session_cache shared:SSL:10m;
add_header Strict-Transport-Security "max-age=31536000" always;
\end{lstlisting}

\subsection{Security Headers}

\begin{lstlisting}[caption=/etc/nginx/snippets/security-headers.conf]
add_header X-Frame-Options "SAMEORIGIN" always;
add_header X-Content-Type-Options "nosniff" always;
add_header X-XSS-Protection "1; mode=block" always;
server_tokens off;
\end{lstlisting}

\subsection{Static Cache}

\begin{lstlisting}[caption=/etc/nginx/snippets/static-cache.conf]
location ~* \.(jpg|jpeg|png|gif|ico|css|js)$ {
    expires 1y;
    add_header Cache-Control "public, immutable";
}
\end{lstlisting}

\subsection{Uso de snippets}

\begin{lstlisting}
server {
    listen 443 ssl http2;
    
    include snippets/ssl-params.conf;
    include snippets/security-headers.conf;
    include snippets/static-cache.conf;
}
\end{lstlisting}

% ===================================
% SECCIÓN 22: EJEMPLO COMPLETO
% ===================================
\section{Ejemplo Práctico Completo}

\subsection{Escenario}

Servidor en \texttt{203.0.113.50} con:
\begin{itemize}
    \item \texttt{empresa.com} → Sitio corporativo
    \item \texttt{blog.empresa.com} → Blog
    \item \texttt{api.empresa.com} → API backend
\end{itemize}

\subsection{Paso 1: Configuración DNS}

\begin{tcolorbox}[colback=yellow!5!white,colframe=yellow!75!black,title=Registros DNS]
\begin{verbatim}
Type  Name                Value           TTL
A     empresa.com         203.0.113.50    3600
A     www.empresa.com     203.0.113.50    3600
A     blog.empresa.com    203.0.113.50    3600
A     api.empresa.com     203.0.113.50    3600
\end{verbatim}
\end{tcolorbox}

\subsection{Paso 2: Crear directorios}

\begin{lstlisting}[language=bash]
sudo mkdir -p /var/www/empresa
sudo mkdir -p /var/www/blog
\end{lstlisting}

\subsection{Paso 3: Configuración del sitio principal}

\begin{lstlisting}[caption=/etc/nginx/sites-available/empresa.conf]
server {
    listen 80;
    server_name empresa.com www.empresa.com;
    
    root /var/www/empresa;
    index index.html;
    
    location / {
        try_files $uri $uri/ /index.html;
    }
}
\end{lstlisting}

\subsection{Paso 4: Configuración del blog}

\begin{lstlisting}[caption=/etc/nginx/sites-available/blog.conf]
server {
    listen 80;
    server_name blog.empresa.com;
    
    root /var/www/blog;
    index index.html;
    
    location / {
        try_files $uri $uri/ =404;
    }
}
\end{lstlisting}

\subsection{Paso 5: Configuración de la API}

\begin{lstlisting}[caption=/etc/nginx/sites-available/api.conf]
server {
    listen 80;
    server_name api.empresa.com;
    
    location / {
        proxy_pass http://localhost:3000;
        proxy_set_header Host $host;
        proxy_set_header X-Real-IP $remote_addr;
    }
}
\end{lstlisting}

\subsection{Paso 6: Habilitar y verificar}

\begin{lstlisting}[language=bash]
# Habilitar sitios
sudo ln -s /etc/nginx/sites-available/empresa.conf \
           /etc/nginx/sites-enabled/
sudo ln -s /etc/nginx/sites-available/blog.conf \
           /etc/nginx/sites-enabled/
sudo ln -s /etc/nginx/sites-available/api.conf \
           /etc/nginx/sites-enabled/

# Verificar configuración
sudo nginx -t

# Recargar Nginx
sudo nginx -s reload

# Probar
curl http://empresa.com
curl http://blog.empresa.com
curl http://api.empresa.com
\end{lstlisting}

% ===================================
% CONCLUSIÓN
% ===================================
\section{Conclusión}

Esta guía cubre los aspectos fundamentales y avanzados de Nginx, desde la configuración básica de virtual hosts hasta optimización de rendimiento y casos de uso complejos.

\subsection{Puntos clave}

\begin{enumerate}
    \item Nginx es un servidor web eficiente y versátil
    \item Los virtual hosts permiten múltiples sitios en un servidor
    \item La configuración es modular y reutilizable
    \item La seguridad debe ser una prioridad
    \item El monitoreo y los logs son esenciales
\end{enumerate}

\subsection{Recursos adicionales}

\begin{itemize}
    \item Documentación oficial: \url{https://nginx.org/en/docs/}
    \item Nginx Admin Guide: \url{https://docs.nginx.com/nginx/admin-guide/}
    \item DigitalOcean Tutorials: \url{https://www.digitalocean.com/community/tags/nginx}
\end{itemize}


\end{document}